\section*{Chapter \ref{ch:g11:marginal-reasoning} --- Marginal Reasoning}
% figdata: figdata/grade-11/01-marginal-reasoning/out/ (life_fit, tax_rates); ledger g11:arkes-blumer-1985; Brenland numbers invented.

\begin{solution}{exo:g11:marginal-reasoning:1}
$26 - 20 = 6$, $30 - 26 = 4$, $36 - 30 = 6$, $46 - 36 = 10$, $60 - 46 = 14$ dalens.
\end{solution}

\begin{solution}{exo:g11:marginal-reasoning:2}
\omterm{def:g11:marginal-reasoning:marginal-benefit}{Marginal benefits}: 30, 20, 12, 6 and 2 points. The third hour brings 12, more than its cost of 10; the fourth brings 6, less. She should revise 3 hours.
\end{solution}

\begin{solution}{exo:g11:marginal-reasoning:3}
$C'(q) = 4 + q$, so $C'(10) = 14$. $C(11) - C(10) = 4 + 0.5 \times (121 - 100) = 14.5$: the step is half a dalen above the slope.
\end{solution}

\begin{solution}{exo:g11:marginal-reasoning:4}
From the curve, about 62.5 years at 100 dollars, 73.7 at 1000 and 85 at 10000: about 11.2 years more for each tenfold rise in spending. The points scatter widely
around the curve.
\end{solution}

\begin{solution}{exo:g11:marginal-reasoning:5}
No: finishing costs 1.5 million dalens for benefits of 1.2 million, a loss of 0.3 million. The 2 million already spent are sunk, so the answer is the same with
3 million.
\end{solution}

\begin{solution}{exo:g11:marginal-reasoning:6}
$N'(q) = 40 - 4q$, zero at $q = 10$, positive before and negative after. $N(10) = 400 - 200 - 50 = 150$.
\end{solution}

\begin{solution}{exo:g11:marginal-reasoning:7}
$B'(q) = 30 - 2q$ and $C'(q) = 6 + q$; they are equal at $q = 8$. $N(8) = (240 - 64) - (48 + 32) = 96$.
\end{solution}

\begin{solution}{exo:g11:marginal-reasoning:8}
Tax $0.20 \times (30000 - 12000) = 3600$ dalens; average rate $3600/30000 = 12\,\%$; marginal rate 20\,\%.
\end{solution}

\begin{solution}{exo:g11:marginal-reasoning:9}
Tax $0.20 \times 28000 + 0.35 \times 60000 + 0.45 \times 20000 = 5600 + 21000 + 9000 = 35600$ dalens; average rate $35600/120000 \approx 29.7\,\%$; marginal rate 45\,\%.
Of an extra \dal{1000} the person keeps \dal{550}.
\end{solution}

\begin{solution}{exo:g11:marginal-reasoning:10}
The \omterm{def:g11:marginal-reasoning:marginal-benefit}{marginal cost} becomes $9 + q$, and $30 - 2q = 9 + q$ gives $q = 7$: one unit fewer. A single charge of 50 does not change any \omterm{def:g11:marginal-reasoning:marginal-quantity}{marginal quantity}: the best
quantity stays 8 and the \omterm{def:g11:marginal-reasoning:net-benefit}{net benefit} falls to 46.
\end{solution}

\begin{solution}{exo:g11:marginal-reasoning:11}
$A(q) = 100/q + q$ and $M(q) = 2q$. At 5: $A = 25$, $M = 10$, the average falls. At 10: $A = M = 20$, the bottom of the average. At 15: $A \approx 21.67$, $M = 30$,
the average rises.
\end{solution}

\begin{solution}{exo:g11:marginal-reasoning:12}
Once the ticket is bought, going to a play costs each buyer the same evening and brings the same pleasure, so the model predicts the same attendance in the three
groups. The first half of the season breaks the prediction: buyers at the normal price used 4.11 tickets against 3.32 and 3.29. The second half fits it: the groups
no longer differed, as if the sunk cost faded from mind after some months. Because the prices were drawn at random, the three groups were alike on average in
their taste for theatre, so the difference can be put down to the price paid, not to the kind of people who pay full price.
\end{solution}

\begin{solution}{exo:g11:marginal-reasoning:13}
The new mean is $(nm + x)/(n + 1) = m + (x - m)/(n + 1)$. It is above $m$ when $x > m$ and below when $x < m$. The added number plays the part of the \omterm{def:g11:marginal-reasoning:marginal-quantity}{marginal
quantity} and the mean that of the average: \cref{prop:g11:marginal-reasoning:average} is the same statement for a smooth total.
\end{solution}

\begin{solution}{exo:g11:marginal-reasoning:14}
$N'(q) = 20.3 - q^2$ is zero at $q = \sqrt{20.3} \approx 4.51$. $N(4) = 81.2 - 21.33 = 59.87$ and $N(5) = 101.5 - 41.67 = 59.83$. The firm should run 4 vans, while
rounding 4.51 gives 5: the \omterm{def:g11:marginal-reasoning:net-benefit}{net benefit} falls faster above its peak than below it, so the nearest whole number is not the best one. With whole units, compare the
whole numbers on both sides of the peak, or add a unit only while $N(q + 1) - N(q) > 0$.
\end{solution}

\begin{solution}{exo:g11:marginal-reasoning:15}
At 0: benefit 6000, no tax, total \dal{6000}. At \dal{10000}: benefit $6000 - 5000 = 1000$, no tax, total \dal{11000}. At \dal{14000}: the benefit ended at \dal{12000}, tax
$0.20 \times 2000 = 400$, total \dal{13600}. At \dal{10000}, one more dalen earned cuts the benefit by 50 centimes: a marginal rate of 50\,\%, above the top tax rate of
45\,\%, for a person with a low income.
\end{solution}

\begin{solution}{pb:g11:marginal-reasoning:1}
Money in dalens a day.

\textbf{1.} $R(q) = pq = 6q - 0.01q^2$. At 100: $p = 5$, $R = 500$.

\textbf{2.} $R'(q) = 6 - 0.02q$, so $R'(100) = 4$. $R(101) - R(100) = 606 - 102.01 - 500 = 3.99$.

\textbf{3.} $C'(q) = 0.9 + 0.01q$, so $C'(100) = 1.9$.

\textbf{4.} $C(100) = 120 + 90 + 50 = 260$, so the average cost is 2.60. It includes the rent, 1.20 a loaf at 100 loaves, which no extra loaf changes, and the cheaper early
loaves; the \omterm{def:g11:marginal-reasoning:marginal-benefit}{marginal cost} counts only the next loaf.

\textbf{5.} Yes: it brings about 4 and costs about 1.9 ($C(101) - C(100) = 1.905$).

\textbf{6.} $6 - 0.02q = 0.9 + 0.01q$ gives $0.03q = 5.1$, $q = 170$.

\textbf{7.} $p = 6 - 1.7 = 4.30$.

\textbf{8.} $R = 731$, $C = 120 + 153 + 144.5 = 417.5$, profit $313.5$.

\textbf{9.} At 169: $R = 1014 - 285.61 = 728.39$, $C = 120 + 152.1 + 142.805 = 414.905$, profit 313.485. At 171: $R = 1026 - 292.41 = 733.59$, $C = 120 + 153.9 + 146.205 = 420.105$,
profit 313.485. Both are below 313.5.

\textbf{10.} Average cost $417.5/170 \approx 2.46$, \omterm{def:g11:marginal-reasoning:marginal-benefit}{marginal cost} $0.9 + 1.7 = 2.60$. The marginal is above the average, so the average is rising.

\textbf{11.} $C(q) = 120 + 1.5q + 0.005q^2$ and $C'(q) = 1.5 + 0.01q$.

\textbf{12.} $6 - 0.02q = 1.5 + 0.01q$ gives $q = 150$.

\textbf{13.} $p = 6 - 1.5 = 4.50$.

\textbf{14.} $R = 675$, $C = 120 + 225 + 112.5 = 457.5$, profit $217.5$.

\textbf{15.} The price rose by 0.20, a third of the cost rise. Buyers bear 0.20 a loaf, a third; the bakery bears the other 0.40, and its profit falls by
$313.5 - 217.5 = 96$ a day; 20 loaves a day are no longer bought.

\textbf{16.} No. The oven's cost is paid whatever the number of loaves; it appears in neither $R'(q)$ nor $C'(q)$.

\textbf{17.} Baking more than 150 loaves lowers the profit, since each extra loaf costs more than it brings in; the oven has been paid for either way, and nothing done now
gets the money back.

\textbf{18.} The rent does not depend on $q$: the best quantity stays 150 and the profit falls by 130, to 87.5.

\textbf{19.} Revenue minus the costs that grow with the loaves is $675 - 337.5 = 337.5$. The bakery should stay open as long as the rent is at most 337.5 dalens a day.

\textbf{20.} Before the flour price rise the bakery bakes 170 loaves a day for a profit of 313.5 dalens; after it, 150 loaves for a profit of 217.5 dalens.
\end{solution}
